%%%%%%%%%%%%%%%%%%%%%%%%%%%%%%%%%%%%%%%%%%%%%%%%%
%%%%              INTRODUCTION               %%%%
%%%%%%%%%%%%%%%%%%%%%%%%%%%%%%%%%%%%%%%%%%%%%%%%%

\section{Introduction}
\label{sec:intro}

% ── MOTIVATION & PROBLEM DOMAIN ──────────────────────────────────────────────
Large language models have substantially advanced biomedical NLP, enabling
question answering, literature summarisation, and clinical note processing at
scale.  A persistent concern, however, is that these models can generate
plausible yet unsupported claims — so-called hallucinations — which pose
serious risks in clinical settings.  The standard mitigation strategy grounds
every LLM output in retrieved evidence through Retrieval-Augmented Generation
(RAG).  TREC BioGen 2024/2025~\cite{biogen2024} formalises this idea as a
shared task: given a biomedical question, retrieve relevant PubMed abstracts,
generate a paragraph-length answer, and cite supporting evidence at the
sentence level.

Several characteristics make this task particularly demanding.
Biomedical queries often mix lay and clinical terminology, whereas PubMed
abstracts use precise scientific notation, creating a vocabulary gap that
hampers lexical matching.  There is also an inherent tension between precision
and recall: missing a key supporting study can be as harmful as retrieving a
contradicting one.  Furthermore, the answer must ground every claim with
citation-level granularity, and evaluation spans both IR metrics (MAP, NDCG)
and generation quality metrics such as citation accuracy.

% ── PROJECT STRUCTURE ─────────────────────────────────────────────────────────
This project is organised into three phases.
\textbf{Phase~1} builds and evaluates a complete
IR pipeline: we index 4{,}194 PubMed abstracts in OpenSearch, implement and
tune five retrieval strategies (BM25, LM-Jelinek-Mercer, LM-Dirichlet, dense
KNN, and Reciprocal Rank Fusion), and evaluate against citation-derived
relevance judgements using NDCG@100, MAP, MRR, P@10, and R@100.
Our best Phase~1 result is \textbf{KNN\,(MedCPT) with NDCG@100\,=\,0.8077,
MAP\,=\,0.6143, R@100\,=\,0.8933} on 33 held-out test queries;
RRF (BM25+MedCPT, $k$=60) reaches NDCG@100\,=\,0.8245.%
\footnote{The full Phase~1 notebook is available at:
\url{https://colab.research.google.com/drive/1UvDhFXYEzZxse6gknDyBSUaq0EAHZkFV?usp=sharing}.}
\textbf{Phase~2} examines the positional embeddings, contextual embeddings, and self-attention as input embeddings go through different layers (with MedCPT Query Encoder and MedCPT Cross Encoder). On top of the Phase 1, it adds cross-encoder reranking of retrieved evidence, LLM-generated paragraph answers with per-sentence citations, and LLM-as-a-Judge evaluation using two different metrics: sentence alignment and answer entailment, aswell as manual analysis of results.
Our best Phase~2 result is \textbf{Gemma 4 31B IT with temperature = 1.0, top-p = 0.7, and a tuned user prompt} on 33 held-out test queries.
\footnote{The full Phase~2 was not developed using Google Colab; the full code can be visualized in the submission documents with title phase2.ipynb}
\textbf{Phase~3} will introduce a ReAct-style agentic loop for multi-step
biomedical synthesis.

% ── CONTRIBUTIONS ─────────────────────────────────────────────────────────────
The implementation builds upon the laboratory exercises of the NLP\&S
course~\cite{opensearch}, adapting their indexing, encoding, and evaluation
patterns to the TREC BioGen domain.

% ── PAPER STRUCTURE ───────────────────────────────────────────────────────────
\noindent
This report is organised as follows.
Section~\ref{sec:biogen} describes the BioGen NL Agent pipeline across three phases.
Section~\ref{sec:eval_setup} details the experimental setup (datasets, metrics, protocols).
Section~\ref{sec:eval_results} presents Phase~1 and Phase~2 results and discussion.
Section~\ref{sec:conclusions} concludes and outlines future work.
